\documentclass[10pt]{article}
\usepackage[a4paper,margin=1in]{geometry}
\usepackage{amsmath,amssymb,amsthm}
\usepackage{booktabs}
\usepackage{array}
\usepackage{enumitem}
\usepackage{microtype}
\usepackage{hyperref}
\hypersetup{colorlinks=true,linkcolor=blue,citecolor=blue,urlcolor=blue}

\newtheorem{lemma}{Lemma}
\newtheorem{proposition}{Proposition}
\newtheorem{remark}{Remark}
\newcommand{\RR}{\mathbb R}
\newcommand{\QQ}{\mathbb Q}
\newcommand{\dd}{\,\mathrm d}
\newcommand{\norm}[1]{\left\lVert #1\right\rVert}
\newcommand{\Gammaq}{\Gamma_{q,r}}

\title{Cross-Prime Gram Coherence in a Literal\\
Twin-Prime Dynamics: A Finite Phase Diagram\\
and Sign-Flip Obstruction}
\author{Liang Wang\\
School of Mathematics and Statistics\\
Huazhong University of Science and Technology (HUST)\\
Wuhan, China}
\date{28 August 2026}

\begin{document}
\maketitle

\begin{abstract}
The preceding finite shell study found a potentially misleading combination:
a four-block scalar attachment can be small while the physical output energy
of the prime shell is amplified.  This paper follows the next structural
question without changing the literal operator or source.  For the physical
prime-component output vectors $g_q$ we form the source-native Gram matrix
$G_{q,r}=\langle g_q,g_r\rangle$ and the normalized squared coherence
$\Gamma_{q,r}=G_{q,r}^2/(G_{q,q}G_{r,r})$.  We prove exactly a conditional
coherence-accumulation inequality: a positive coherence floor and diagonal
balance force a lower bound on the aggregate energy ratio.  An exact-rational
finite scan of 18 rows and 1,380 unordered pairs then separates two regimes.
Seventeen rows are pairwise positive, but the early $(N,H,Q,z,s)=(256,38,27,5,1)$
row contains three negative cross-prime pairs, one with squared coherence about
$1.37\times 10^{-7}$.  Conversely, eight late-shell rows satisfy the finite
thresholds $\Gamma_{q,r}\ge 9/25$ and $d_{\min}/d_{\max}\ge4/5$ for every pair,
while all 18 rows have aggregate energy ratio greater than one.  The result is
a finite sign/coherence phase diagram and a precise obstruction to promoting
late-shell behavior to a uniform growing-shell theorem; it is not an
asymptotic twin-prime result or an arithmetic $L^2$ estimate.
\end{abstract}

\section{Position on the route map}

The working line studies a single literal dynamical family intended to model a
prime-shell contribution to the twin-prime problem.  Earlier releases split
the finite physical operator into prime components and documented scalar
cancellation.  The immediately preceding release, TPC-288, retained the full
output vectors and proved the exact Gram identity while certifying full rank on
a finite growth/control grid \cite{tpc288}.  Its central warning was that a
small scalar attachment is not an $L^2$ saving.

TPC-289 asks the smallest follow-up question suggested by that warning:
whether the cross-prime terms themselves have a stable sign and size.  This is
more informative than another rank computation, but it remains deliberately
finite.  The distinction between exact identities, finite numerical
certificates, and open asymptotic gates is maintained throughout.

\section{The frozen physical model}

Let $I=I_N$ be the integer interval used by the frozen source engine and let
$\beta\in\QQ^I$ be its source vector.  For an odd prime $q$, define the
centered residue kernel
\[
 B_q(u,t)=\mathbf 1_{q\nmid u}\mathbf 1_{q\nmid t}
 \left(\mathbf 1_{u\equiv t\pmod q}-\frac{1}{q-1}\right).
\]
With
\[
 K_{H,s}(h)=\frac{H^{2s}}{(H^2+h^2)^s},
 \qquad
 D_q(u,t)=\mathbf 1_{u\ne t}B_q(u,t),
\]
the literal deleted-diagonal component is
\[
 (A_q)_{u,t}=qK_{H,s}(u-t)D_q(u,t),
 \qquad g_q=A_q\beta.
\]
For a finite shell $S$ write
\[
 g_S=\sum_{q\in S}g_q,\qquad
 G_{q,r}=\langle g_q,g_r\rangle,\qquad d_q=G_{q,q}=\norm{g_q}_2^2.
\]
No four-block scalar attachment is applied in this paper.  The source, kernel,
deleted diagonal, and prime shell are exactly those inherited from TPC-288.

For $q\ne r$ with $d_qd_r>0$, define
\[
 \Gamma_{q,r}=\frac{G_{q,r}^2}{d_qd_r}.
\]
The square records the size of a cross term independently of its sign; the
sign of $G_{q,r}$ is retained as a separate datum.  Expanding the aggregate
square gives the exact energy decomposition
\begin{equation}
 \frac{\norm{g_S}_2^2}{\sum_{q\in S}d_q}
 =1+\frac{\sum_{q\ne r}G_{q,r}}{\sum_{q\in S}d_q},
 \label{eq:energy}
\end{equation}
where the off-diagonal sum is over ordered pairs.

\section{Exact coherence structure}

\begin{lemma}[Gram positivity and coherence bound]
For every finite shell, $G$ is positive semidefinite.  If $d_q,d_r>0$, then
\begin{equation}
 0\leq \Gamma_{q,r}\leq 1.
 \label{eq:cauchy}
\end{equation}
\end{lemma}

\begin{proof}
For any real coefficient vector $c$,
\[
 c^T Gc=\sum_{q,r}c_qc_r\langle g_q,g_r\rangle
 =\norm{\sum_qc_qg_q}_2^2\geq0.
\]
The lower bound in \eqref{eq:cauchy} is a square and the upper bound is
Cauchy--Schwarz divided by $d_qd_r>0$.
\end{proof}

\begin{proposition}[Conditional coherence accumulation]
Let $k=|S|$, $d_- =\min_{q\in S}d_q$, and
$d_+=\max_{q\in S}d_q$.  Suppose that for some $0\leq\eta,\delta\leq1$,
\[
 G_{q,r}>0,\qquad \Gamma_{q,r}\geq\eta^2\quad(q\ne r),
 \qquad \frac{d_-}{d_+}\geq\delta.
 \label{eq:hypotheses}
\]
Then
\begin{equation}
 \frac{\norm{g_S}_2^2}{\sum_qd_q}
 \geq 1+\eta\delta(k-1).
 \label{eq:accumulation}
\end{equation}
\end{proposition}

\begin{proof}
The positivity assumption selects the positive square-root branch in the
definition of $\Gamma_{q,r}$, so
\[
 G_{q,r}\geq\eta\sqrt{d_qd_r}\geq\eta d_-.
\]
There are $k(k-1)$ ordered off-diagonal terms.  Since
$\sum_qd_q\leq kd_+$, equation \eqref{eq:energy} yields
\[
 R_E\geq1+\frac{k(k-1)\eta d_-}{kd_+}
 \geq1+\eta\delta(k-1).
\]
\end{proof}

\begin{remark}
The proposition is a finite vector lemma, not a claim that its hypotheses
hold for all prime shells.  If even one cross term is negative, its hypotheses
fail, irrespective of whether the total energy ratio happens to exceed one.
This logical separation is the obstruction tested below.
\end{remark}

\section{Finite protocol}

The certificate uses exact Python `Fraction` arithmetic and the frozen
TPC-268 engine.  The 18 rows consist of eight $s=2$ growth anchors, four
selected $s=1$ exponent-crossover rows, and six $s=2$ source-control rows at
$Q=70$.  The shell anchors range from three to 17 primes.  The six controls
use $H\in\{48,52\}$ and $z\in\{3,5,7\}$.

For every row the producer reconstructs each $g_q$, forms every rational Gram
entry, checks \eqref{eq:cauchy}, records the sign census, and computes the
exact energy ratio in \eqref{eq:energy}.  The finite ``strong block'' test uses
\[
 \eta=\frac35,\qquad \delta=\frac45,
 \qquad \eta^2=\frac9{25}.
\]
Thus a 15-prime row passing the test has the conditional lower envelope
$R_E\geq1+(3/5)(4/5)14=193/25$, while a 17-prime row has
$R_E\geq217/25$.

The producer locks the TPC-288 source code and result, as well as the frozen
engine, by normalized-LF SHA-256.  An independent checker reverses the output
summation order and reconstructs the complete canonical JSON document.  A
separate stress script rejects ten mutations, including altered signs,
thresholds, row counts, energy values, and provenance.  Normal and optimized
Python executions are required to have byte-identical output.

\section{Results: a finite sign/coherence phase diagram}

Table~\ref{tab:representative} gives representative rows.  ``min $\Gamma$''
is the smallest squared coherence over unordered pairs; ``$+$/$-$'' is the
pair sign census; and ``strong'' means that all three hypotheses in
\eqref{eq:hypotheses} pass with the declared $(\eta,\delta)$.

\begin{table}[ht]
\centering
\scriptsize
\begin{tabular}{@{}lrrrrrrrrr@{}}
\toprule
axis & $N$ & $H$ & $Q$ & $s$ & $k$ & min $\Gamma$ & $+/-$ & $R_E$ & strong\\
\midrule
growth &128&24&9 &2&3 &0.292716&3/0&2.124905&no\\
growth &192&32&16&2&5 &0.251582&10/0&3.443701&no\\
growth &256&38&27&2&7 &0.000949&21/0&2.622578&no\\
growth &384&50&40&2&10&0.207063&45/0&7.106564&no\\
growth &512&58&60&2&13&0.331366&78/0&10.992677&no\\
growth &512&58&70&2&15&0.444464&105/0&13.558338&yes\\
growth &512&58&90&2&17&0.806422&136/0&16.439313&yes\\
cross &256&38&27&1&7&$1.37\!\times10^{-7}$&18/3&2.166576&no\\
cross &384&50&40&1&10&0.032697&45/0&4.557433&no\\
cross &512&58&70&1&15&0.058904&105/0&8.900964&no\\
control&384&48&70&2&15&0.721183&105/0&14.322934&yes\\
control&384&52&70&2&15&0.642607&105/0&14.107249&yes\\
\bottomrule
\end{tabular}
\caption{Selected exact-rational rows; decimals are display strings only.}
\label{tab:representative}
\end{table}

Across all 18 rows there are 1,380 unordered pair comparisons.  Seventeen
rows are pairwise positive and one row has three negative pairs; no pair is
zero.  Every row has $R_E>1$.  The exceptional row is
$(256,38,27,5,1)$, with shell
$\{29,31,37,41,43,47,53\}$.  Its negative pairs are
\[
 (29,53),\qquad(31,53),\qquad(41,53).
\]
The squared coherence of $(31,53)$ is approximately
$1.3745664\times10^{-7}$, while $(29,53)$ has negative Gram cross term and
squared coherence approximately $0.0105163$.  Thus the failure is not merely
a rounding ambiguity: the signs are exact rational signs and one pair is
nearly orthogonal.

The late block consists of eight rows: the $Q=70$ and $Q=90$ $s=2$ growth
anchors, together with the six $Q=70$ controls.  Every pair is positive,
every squared coherence is at least $9/25$, and the diagonal ratio is at least
$4/5$.  The conditional proposition therefore applies row by row.  The
finite certificate also finds two exact groups in which the three cutoff
controls at a fixed height have the same recorded coherence/energy signature;
this is a warning that these finite controls should not be counted as an
asymptotically independent source family.

\section{What is and is not established}

The positive late block is useful because it supplies a quantitative mechanism
for the energy amplification seen in TPC-288: coherent cross terms accumulate
linearly in the number of components.  The early sign flip is equally useful:
it refutes the tempting uniform rule
\[
 \text{``all declared shell cross terms are positive and have a fixed
 coherence floor.''}
\]
The observed total $R_E>1$ in the exceptional row shows that aggregate
amplification does not require pairwise positivity either.

The result is therefore an obstruction paper on the route map, not a proof or
disproof of the twin-prime conjecture.  In particular, we do not claim a
growing-shell bound, a uniform theorem over literal sources, an arithmetic
$L^2$ estimate, a fixed-power credit, or a Gate-B conclusion.  The next
minimal questions are whether an adaptive shell weighting can exploit the
coherence without violating the physical operator, and whether a restricted
source class admits a uniform coherence estimate.

\section{Reproducibility and claim firewall}

The complete artifact is in
\href{../README.md}{the project README}.  The canonical result is
`results/tpc289\_certificate.json`; the proof and derivation packages state the
exact scope.  The local fallback evaluator reports:

\begin{verbatim}
EXACT: Gram PSD, 0 <= Gamma <= 1, conditional accumulation envelope
FINITE: 17/18 positive rows; 3 negative pairs in one row
FINITE: 8-row eta=3/5, delta=4/5 strong block; 18/18 R_E > 1
OPEN: source-restricted/growing-shell coherence and arithmetic L2
NONE: twin-prime conclusion
\end{verbatim}

The Session-named Route-A/Route-B evaluator files are not present in this
checkout, so no official evaluator pass is claimed.  The independent replay,
mutation audit, and bridge checker are the available fail-closed evidence.

\bibliographystyle{plain}
\bibliography{references}

\end{document}
